% ============================================================
% SECTION 6: APPLICATION DEVELOPMENT
% ============================================================
\section{Phát triển Ứng dụng}
\label{sec:application-development}

% --- 6.1 Backend API ---
\subsection{Backend API Server (Golang)}
\label{sec:backend-api}

Backend API được xây dựng bằng \textbf{Golang} với framework \textbf{Gin} cho REST endpoints và package \texttt{gorilla/websocket} cho WebSocket hub.

\subsubsection{REST API Endpoints}
\bigskip
\begin{table}[H]
\centering
\caption[Các REST API Endpoints chính]{Danh sách REST API Endpoints}
\label{tab:rest-api}
\begin{tabular}{|p{3cm}|p{2.5cm}|p{6cm}|}
\hline
\textbf{Endpoint} & \textbf{Method} & \textbf{Mô tả} \\
\hline
\texttt{/api/orders} & POST & Tạo đơn hàng mới (Customer) \\
\hline
\texttt{/api/orders/:id} & GET & Lấy thông tin đơn hàng \\
\hline
\texttt{/api/orders/:id/status} & PUT & Cập nhật trạng thái đơn hàng (Driver) \\
\hline
\texttt{/api/trips/:id} & GET & Lấy tóm tắt chuyến đi \\
\hline
\texttt{/api/trips/:id/route} & GET & Lấy full GPS trace cho playback \\
\hline
\texttt{/api/drivers/:id/analytics} & GET & Lấy thống kê hiệu suất tuần \\
\hline
\texttt{/api/drivers/:id/alerts} & GET & Lấy lịch sử cảnh báo \\
\hline
\texttt{/api/admin/heatmap} & GET & Lấy dữ liệu density cho heatmap \\
\hline
\end{tabular}
\end{table}

\subsubsection{WebSocket Hub — Real-Time Push}
\bigskip
\begin{lstlisting}[language=Go, caption={WebSocket Hub — \texttt{hub.go}: quản lý kết nối và broadcast real-time}, label={lst:websocket-hub}]
package websocket

import (
    "strings"
    "time"
    "github.com/gorilla/websocket"
)

type Hub struct {
    clients    map[*Client]bool       // all connected clients
    driverSubs map[string][]*Client  // driver_id → subscribed clients
    register   chan *Client
    unregister chan *Client
    broadcast  chan []byte
}

type Client struct {
    hub      *Hub
    conn     *websocket.Conn
    send     chan []byte
    driverID string
}

func (h *Hub) Run() {
    for {
        select {
        case client := <-h.register:
            h.clients[client] = true
            if clients, ok := h.driverSubs[client.driverID]; ok {
                h.driverSubs[client.driverID] = append(clients, client)
            } else {
                h.driverSubs[client.driverID] = []*Client{client}
            }

        case client := <-h.unregister:
            if _, ok := h.clients[client]; ok {
                delete(h.clients, client)
                delete(h.driverSubs, client.driverID)
                close(client.send)
            }

        case message := <-h.broadcast:
            // message format: "driverID:JSONpayload"
            parts := strings.SplitN(string(message), ":", 2)
            if len(parts) < 2 { continue }
            driverID, payload := parts[0], parts[1]
            for _, client := range h.driverSubs[driverID] {
                select {
                case client.send <- []byte(payload):
                default:
                    close(client.send)
                    delete(h.driverSubs, driverID)
                }
            }
        }
    }
}

// StartKafkaConsumer -- reads from Kafka and pushes to hub
func (h *Hub) StartKafkaConsumer(brokers []string) {
    consumer, _ := kafka.NewConsumer(kafka.ConfigMap{
        "bootstrap.servers":  brokers[0],
        "group.id":          "websocket-consumer",
        "auto.offset.reset":  "latest",
    })
    defer consumer.Close()
    consumer.SubscribeTopics([]string{"processed-updates"}, nil)

    for {
        msg, err := consumer.ReadMessage(5 * time.Second)
        if err == nil {
            h.broadcast <- msg.Value
        }
    }
}
\end{lstlisting}

\bigskip
\begin{lstlisting}[language=Go, caption={Trip Completion — Tính cước phí từ GPS trace (\texttt{internal/cassandra/trips.go})}, label={lst:trip-completion}]
func (c *Client) CompleteTrip(tripID string) error {
    // 1. Get all GPS points for the trip
    points, err := c.GetTripRoute(tripID)
    if err != nil || len(points) == 0 {
        return errors.New("no GPS points for trip")
    }

    // 2. Calculate statistics
    startTime     := points[0].Timestamp
    endTime       := points[len(points)-1].Timestamp
    durationSec   := int(endTime.Sub(startTime).Seconds())
    totalDistance := CalculateTripDistance(points)  // Haversine sum

    var maxSpeed, totalSpeed float64
    speedingViolations := 0
    for _, p := range points {
        if p.Speed > maxSpeed { maxSpeed = p.Speed }
        if p.Speed > 60       { speedingViolations++ }
        totalSpeed += p.Speed
    }
    avgSpeed := totalSpeed / float64(len(points))

    // 3. Pricing: cost = 3.00 + (0.50 * km) + (0.10 * min)
    const baseFare     = 3.00
    const distRate    = 0.50  // per km
    const timeRate    = 0.10  // per minute
    durationMin := float64(durationSec) / 60.0
    tripCost := baseFare + (distRate * totalDistance) + (timeRate * durationMin)
    tripCost = math.Round(tripCost*100) / 100  // round to 2dp

    // 4. Update trip_metadata
    session.Query(`UPDATE trip_metadata SET
        end_time=?, total_distance=?, total_duration=?,
        average_speed=?, max_speed=?, speeding_violations=?,
        trip_cost=?, status='COMPLETED' WHERE trip_id=?`,
        endTime, totalDistance, durationSec, avgSpeed,
        maxSpeed, speedingViolations, tripCost, tripID).Exec()

    log.Printf("Trip %s: cost=%.2f, dist=%.2fkm, dur=%ds",
                tripID, tripCost, totalDistance, durationSec)
    return nil
}
\end{lstlisting}

\bigskip
\noindent
\textbf{Hoạt động của WebSocket Hub:}
\begin{itemize}
    \item Client kết nối WebSocket, gửi: \texttt{\{"action": "subscribe", "driver\_id": "D001"\}}
    \item Hub đăng ký client vào map \texttt{driverSubs["D001"]}.
    \item Kafka consumer đọc message từ \texttt{processed-updates} liên tục.
    \item Consumer gửi đến hub qua channel \texttt{broadcast}.
    \item Hub tra cứu \texttt{driverSubs} và push đến tất cả client đã subscribe.
    \item Browser nhận JSON và cập nhật vị trí marker trên bản đồ.
\end{itemize}

% --- 6.2 Frontend Interface ---
\subsection{Giao diện người dùng (React + Leaflet)}
\label{sec:frontend}

Frontend được xây dựng bằng \textbf{React} (Vite) kết hợp \textbf{Leaflet} để hiển thị bản đồ OpenStreetMap.

\bigskip
\begin{figure}[H]
    \centering
    \includegraphics[width=0.85\textwidth]{images/frontend-customer.png}
    \caption[Giao diện khách hàng — Theo dõi tài xế real-time]{Giao diện Customer Dashboard — Bản đồ theo dõi vị trí tài xế real-time với marker di chuyển mỗi giây}
    \label{fig:frontend-customer}
\end{figure}

\bigskip
\begin{figure}[H]
    \centering
    \includegraphics[width=0.85\textwidth]{images/frontend-driver.png}
    \caption[Giao diện tài xế — Trạng thái đơn hàng]{Giao diện Driver App — Cập nhật trạng thái giao hàng: PENDING $\rightarrow$ ACCEPTED $\rightarrow$ IN\_TRANSIT $\rightarrow$ DELIVERED}
    \label{fig:frontend-driver}
\end{figure}

\bigskip
\begin{figure}[H]
    \centering
    \includegraphics[width=0.85\textwidth]{images/frontend-admin.png}
    \caption[Giao diện quản trị — Dashboard]{Giao diện Admin Dashboard — Analytics charts, alert feed, heatmap overlay}
    \label{fig:frontend-admin}
\end{figure}

\bigskip
\noindent
\textbf{Các tính năng chính của frontend:}

\bigskip
\begin{itemize}
    \item \textbf{Customer View (Hình \ref{fig:frontend-customer}):}
    \begin{itemize}
        \item Nút "Place Order" tạo đơn hàng qua REST API.
        \item Leaflet Map hiển thị marker tài xế, cập nhật vị trí mỗi giây qua WebSocket.
        \item ETA countdown hiển thị thời gian đến ước tính.
        \item Toast notification khi tài xế cách < 500m (proximity alert).
    \end{itemize}

    \item \textbf{Driver View (Hình \ref{fig:frontend-driver}):}
    \begin{itemize}
        \item Danh sách đơn hàng đang chờ (PENDING).
        \item Nút "Accept / Reject" -- cập nhật trạng thái qua REST.
        \item Trạng thái: \texttt{ACCEPTED -\textgreater{} PICKING\_UP -\textgreater{} IN\_TRANSIT -\textgreater{} ARRIVING -\textgreater{} DELIVERED}.
        \item Khi DELIVERED -\textgreater{} gọi \texttt{CompleteTrip()} -\textgreater{} tính cước phí.
    \end{itemize}

    \item \textbf{Admin View (Hình \ref{fig:frontend-admin}):}
    \begin{itemize}
        \item Charts (Recharts) hiển thị hiệu suất tài xế theo tuần.
        \item Alert feed real-time từ WebSocket.
        \item Heatmap overlay (Leaflet.heat) hiển thị khu vực nhu cầu cao.
        \item Trip playback -- animation GPS trace từ API \texttt{GET /api/trips/:id/route}.
    \end{itemize}
\end{itemize}

\bigskip
\noindent
\textbf{WebSocket Client Implementation (React Hook):}
\bigskip
\begin{lstlisting}[language=JavaScript, caption={React Hook — \texttt{useDriverLocation.js}: WebSocket subscription cho real-time location}, label={lst:websocket-client}]
import { useState, useEffect } from "react";

export function useDriverLocation(driverId) {
  const [location, setLocation] = useState(null);

  useEffect(() => {
    const protocol = window.location.protocol === "https:" ? "wss:" : "ws:";
    const ws = new WebSocket(`${protocol}//localhost:8080/ws`);

    ws.onopen = () => {
      ws.send(JSON.stringify({
        action:    "subscribe",
        driver_id: driverId,
      }));
    };

    ws.onmessage = (event) => {
      const data = JSON.parse(event.data);

      if (data.type === "location_update") {
        setLocation({
          lat:   data.latitude,
          lng:   data.longitude,
          speed: data.speed,
          eta:   data.eta_seconds,
        });
      }

      if (data.type === "alert") {
        showToast(data.message); // Proximity / Speeding notification
      }
    };

    ws.onerror = (err) => console.error("WebSocket error:", err);
    return () => ws.close();
  }, [driverId]);

  return location;
}
\end{lstlisting}
